\section{The splitting principle}
\label{section-splitting-principle}

\noindent
In our setting it is not so easy to say what the splitting principle
exactly says/is. Here is a possible formulation.

\begin{lemma}
\label{lemma-splitting-principle}
Let $(S, \delta)$ be as in Situation \ref{situation-setup}. Let $X$ be locally
of finite type over $S$. Let $\mathcal{E}_i$ be a finite collection of
locally free $\mathcal{O}_X$-modules of rank $r_i$. There exists a projective
flat morphism $\pi : P \to X$ of relative dimension $d$ such that
\begin{enumerate}
\item for any morphism $f : Y \to X$ the map
$\pi_Y^* : \CH_*(Y) \to \CH_{* + d}(Y \times_X P)$ is injective, and
\item each $\pi^*\mathcal{E}_i$ has a filtration
whose successive quotients $\mathcal{L}_{i, 1}, \ldots, \mathcal{L}_{i, r_i}$
are invertible ${\mathcal O}_P$-modules.
\end{enumerate}
Moreover, when (1) holds the restriction map $A^*(X) \to A^*(P)$
(Remark \ref{remark-pullback-cohomology}) is injective.
\end{lemma}

\begin{proof}
We may assume $r_i \geq 1$ for all $i$. We will prove the lemma by induction
on $\sum (r_i - 1)$. If this integer is $0$, then $\mathcal{E}_i$
is invertible for all $i$ and we conclude by taking $\pi = \text{id}_X$.
If not, then we can pick an $i$ such that $r_i > 1$ and consider the
morphism $\pi_i : P_i = \mathbf{P}(\mathcal{E}_i) \to X$.
We have a short exact sequence
$$
0 \to \mathcal{F} \to \pi_i^*\mathcal{E}_i \to \mathcal{O}_{P_i}(1) \to 0
$$
of finite locally free $\mathcal{O}_{P_i}$-modules of ranks $r_i - 1$,
$r_i$, and $1$. Observe that $\pi_i^*$ is injective on chow groups
after any base change by the projective bundle formula
(Lemma \ref{lemma-chow-ring-projective-bundle}).
By the induction hypothesis applied to the finite locally free
$\mathcal{O}_{P_i}$-modules $\mathcal{F}$ and $\pi_{i'}^*\mathcal{E}_{i'}$
for $i' \not = i$, we find a morphism $\pi : P \to P_i$ with
properties stated as in the lemma. Then the composition
$\pi_i \circ \pi : P \to X$ does the job. Some details omitted.
\end{proof}

\begin{remark}
\label{remark-the-proof-shows-more}
The proof of Lemma \ref{lemma-splitting-principle}
shows that the morphism $\pi : P \to X$ has the following additional
properties:
\begin{enumerate}
\item $\pi$ is a finite composition of projective space bundles
associated to locally free modules of finite constant rank, and
\item for every $\alpha \in \CH_k(X)$ we have
$\alpha = \pi_*(\xi_1 \cap \ldots \cap \xi_d \cap \pi^*\alpha)$
where $\xi_i$ is the first Chern class of some invertible
$\mathcal{O}_P$-module.
\end{enumerate}
The second observation follows from the first and
Lemma \ref{lemma-cap-projective-bundle}.
We will add more observations here as needed.
\end{remark}

\noindent
Let $(S, \delta)$, $X$, and $\mathcal{E}_i$ be as in
Lemma \ref{lemma-splitting-principle}.
The {\it splitting principle} refers to the practice of symbolically writing
$$
c(\mathcal{E}_i) = \prod (1 + x_{i, j})
$$
The symbols $x_{i, 1}, \ldots, x_{i, r_i}$ are called the {\it Chern roots}
of $\mathcal{E}_i$. In other words, the $p$th Chern class of $\mathcal{E}_i$
is the $p$th elementary symmetric function in the Chern roots.
The usefulness of the splitting principle comes from the assertion that
in order to prove a polynomial relation among Chern classes of the
$\mathcal{E}_i$ it is enough to prove the corresponding relation among the
Chern roots.

\medskip\noindent
Namely, let $\pi : P \to X$ be as in Lemma \ref{lemma-splitting-principle}.
Recall that there is a canonical $\mathbf{Z}$-algebra map
$\pi^* : A^*(X) \to A^*(P)$, see
Remark \ref{remark-pullback-cohomology}. The injectivity of $\pi_Y^*$
on Chow groups for every $Y$ over $X$, implies that the map
$\pi^* : A^*(X) \to A^*(P)$ is injective (details omitted).
We have
$$
\pi^*c(\mathcal{E}_i) = \prod (1 + c_1(\mathcal{L}_{i, j}))
$$
by Lemma \ref{lemma-chern-filter-by-linebundles}. Thus we may think of the
Chern roots $x_{i, j}$ as the elements $c_1(\mathcal{L}_{i, j}) \in A^*(P)$
and the displayed equation as taking place in $A^*(P)$ after
applying the injective map $\pi^* : A^*(X) \to A^*(P)$ to the left
hand side of the equation.

\medskip\noindent
To see how this works, it is best to give some examples.

\begin{lemma}
\label{lemma-chern-classes-dual}
In Situation \ref{situation-setup} let $X$ be locally of finite type over $S$.
Let $\mathcal{E}$ be a finite locally free $\mathcal{O}_X$-module
with dual $\mathcal{E}^\vee$. Then
$$
c_i(\mathcal{E}^\vee) = (-1)^i c_i(\mathcal{E})
$$
in $A^i(X)$.
\end{lemma}

\begin{proof}
Choose a morphism $\pi : P \to X$ as in
Lemma \ref{lemma-splitting-principle}.
By the injectivity of $\pi^*$ (after any base change)
it suffices to prove the relation between
the Chern classes of $\mathcal{E}$ and $\mathcal{E}^\vee$
after pulling back to $P$. Thus we may assume there
exist invertible $\mathcal{O}_X$-modules
${\mathcal L}_i$, $i = 1, \ldots, r$
and a filtration
$$
0 = \mathcal{E}_0 \subset \mathcal{E}_1 \subset \mathcal{E}_2
\subset \ldots \subset \mathcal{E}_r = \mathcal{E}
$$
such that $\mathcal{E}_i/\mathcal{E}_{i - 1} \cong \mathcal{L}_i$.
Then we obtain the dual filtration
$$
0 = \mathcal{E}_r^\perp \subset \mathcal{E}_1^\perp \subset \mathcal{E}_2^\perp
\subset \ldots \subset \mathcal{E}_0^\perp = \mathcal{E}^\vee
$$
such that $\mathcal{E}_{i - 1}^\perp/\mathcal{E}_i^\perp \cong
\mathcal{L}_i^{\otimes -1}$.
Set $x_i = c_1(\mathcal{L}_i)$.
Then $c_1(\mathcal{L}_i^{\otimes -1}) = - x_i$
by Lemma \ref{lemma-c1-cap-additive}.
By Lemma \ref{lemma-chern-filter-by-linebundles}
we have
$$
c(\mathcal{E}) = \prod\nolimits_{i = 1}^r (1 + x_i)
\quad\text{and}\quad
c(\mathcal{E}^\vee) = \prod\nolimits_{i = 1}^r (1 - x_i)
$$
in $A^*(X)$. The result follows from a formal computation
which we omit.
\end{proof}

\begin{lemma}
\label{lemma-chern-classes-tensor-product}
In Situation \ref{situation-setup} let $X$ be locally of finite type over $S$.
Let $\mathcal{E}$ and $\mathcal{F}$ be a finite locally free
$\mathcal{O}_X$-modules of ranks $r$ and $s$. Then we have
$$
c_1(\mathcal{E} \otimes \mathcal{F})
=
r c_1(\mathcal{F}) + s c_1(\mathcal{E})
$$
$$
c_2(\mathcal{E} \otimes \mathcal{F})
=
r c_2(\mathcal{F}) + s c_2(\mathcal{E}) +
{r \choose 2} c_1(\mathcal{F})^2 +
(rs - 1) c_1(\mathcal{F})c_1(\mathcal{E}) +
{s \choose 2} c_1(\mathcal{E})^2
$$
and so on in $A^*(X)$.
\end{lemma}

\begin{proof}
Arguing exactly as in the proof of Lemma \ref{lemma-chern-classes-dual}
we may assume we have
invertible $\mathcal{O}_X$-modules
${\mathcal L}_i$, $i = 1, \ldots, r$
${\mathcal N}_i$, $i = 1, \ldots, s$
filtrations
$$
0 = \mathcal{E}_0 \subset \mathcal{E}_1 \subset \mathcal{E}_2
\subset \ldots \subset \mathcal{E}_r = \mathcal{E}
\quad\text{and}\quad
0 = \mathcal{F}_0 \subset \mathcal{F}_1 \subset \mathcal{F}_2
\subset \ldots \subset \mathcal{F}_s = \mathcal{F}
$$
such that $\mathcal{E}_i/\mathcal{E}_{i - 1} \cong \mathcal{L}_i$
and such that $\mathcal{F}_j/\mathcal{F}_{j - 1} \cong \mathcal{N}_j$.
Ordering pairs $(i, j)$ lexicographically
we obtain a filtration
$$
0 \subset \ldots \subset
\mathcal{E}_i \otimes \mathcal{F}_j
+
\mathcal{E}_{i - 1} \otimes \mathcal{F}
\subset \ldots \subset \mathcal{E} \otimes \mathcal{F}
$$
with successive quotients
$$
\mathcal{L}_1 \otimes \mathcal{N}_1,
\mathcal{L}_1 \otimes \mathcal{N}_2,
\ldots,
\mathcal{L}_1 \otimes \mathcal{N}_s,
\mathcal{L}_2 \otimes \mathcal{N}_1,
\ldots,
\mathcal{L}_r \otimes \mathcal{N}_s
$$
By Lemma \ref{lemma-chern-filter-by-linebundles}
we have
$$
c(\mathcal{E}) = \prod (1 + x_i),
\quad
c(\mathcal{F}) = \prod (1 + y_j),
\quad\text{and}\quad
c(\mathcal{E} \otimes \mathcal{F}) = \prod (1 + x_i + y_j),
$$
in $A^*(X)$. The result follows from a formal computation
which we omit.
\end{proof}

\begin{remark}
\label{remark-equalities-nonconstant-rank}
The equalities proven above remain true even when we work with
finite locally free
$\mathcal{O}_X$-modules whose rank is allowed to be nonconstant.
In fact, we can work with polynomials in the rank and the
Chern classes as follows. Consider the  graded polynomial ring
$\mathbf{Z}[r, c_1, c_2, c_3, \ldots]$
where $r$ has degree $0$ and $c_i$ has degree $i$. Let
$$
P \in \mathbf{Z}[r, c_1, c_2, c_3, \ldots]
$$
be a homogeneous polynomial of degree $p$. Then for any finite locally
free $\mathcal{O}_X$-module $\mathcal{E}$ on $X$ we can consider
$$
P(\mathcal{E}) =
P(r(\mathcal{E}), c_1(\mathcal{E}), c_2(\mathcal{E}), c_3(\mathcal{E}), \ldots)
\in A^p(X)
$$
see Remark \ref{remark-extend-to-finite-locally-free} for notation and
conventions. To prove relations among these polynomials (for multiple
finite locally free modules) we can work locally on $X$ and use the splitting
principle as above. For example, we claim that
$$
c_2(\SheafHom_{\mathcal{O}_X}(\mathcal{E}, \mathcal{E})) =
P(\mathcal{E})
$$
where $P = 2rc_2 - (r - 1)c_1^2$.
Namely, since $\SheafHom_{\mathcal{O}_X}(\mathcal{E}, \mathcal{E}) =
\mathcal{E} \otimes \mathcal{E}^\vee$ this follows easily from
Lemmas \ref{lemma-chern-classes-dual} and
\ref{lemma-chern-classes-tensor-product}
above by decomposing $X$ into parts where the rank
of $\mathcal{E}$ is constant as in
Remark \ref{remark-extend-to-finite-locally-free}.
\end{remark}

\begin{example}
\label{example-power-sum}
For every $p \geq 1$ there is a unique homogeneous polynomial
$P_p \in \mathbf{Z}[c_1, c_2, c_3, \ldots]$ of degree $p$
such that, for any $n \geq p$ we have
$$
P_p(s_1, s_2, \ldots, s_p) = \sum x_i^p
$$
in $\mathbf{Z}[x_1, \ldots, x_n]$ where $s_1, \ldots, s_p$ are the
elementary symmetric polynomials in $x_1, \ldots, x_n$, so
$$
s_i = \sum\nolimits_{1 \leq j_1 < \ldots < j_i \leq n}
x_{j_1}x_{j_2} \ldots x_{j_i}
$$
The existence of $P_p$ comes from the well known fact that
the elementary symmetric functions generate the ring of
all symmetric functions over the integers. Another way to
characterize $P_p \in \mathbf{Z}[c_1, c_2, c_3, \ldots]$ is that we have
$$
\log(1 + c_1 + c_2 + c_3 + \ldots) =
\sum\nolimits_{p \geq 1} (-1)^{p - 1}\frac{P_p}{p}
$$
as formal power series. This is clear by writing
$1 + c_1 + c_2 + \ldots = \prod (1 + x_i)$ and applying
the power series for the logarithm function. Expanding the left
hand side we get
\begin{align*}
& (c_1 + c_2 + \ldots) - (1/2)(c_1 + c_2 + \ldots)^2 +
(1/3)(c_1 + c_2 + \ldots)^3 - \ldots \\
& =
c_1 + (c_2 - (1/2)c_1^2) + (c_3 - c_1c_2 + (1/3)c_1^3) + \ldots
\end{align*}
In this way we find that
\begin{align*}
P_1 & = c_1, \\
P_2 & = c_1^2 - 2c_2, \\
P_3 & = c_1^3 - 3c_1c_2 + 3c_3, \\
P_4 & = c_1^4 - 4c_1^2c_2 + 4c_1c_3 + 2c_2^2 - 4c_4,
\end{align*}
and so on. Since the Chern classes of a finite locally free
$\mathcal{O}_X$-module $\mathcal{E}$ are the elementary symmetric
polynomials in the Chern roots $x_i$, we see that
$$
P_p(\mathcal{E}) = \sum x_i^p
$$
For convenience we set $P_0 = r$ in $\mathbf{Z}[r, c_1, c_2, c_3, \ldots]$
so that $P_0(\mathcal{E}) = r(\mathcal{E})$ as a bivariant class
(as in Remarks \ref{remark-extend-to-finite-locally-free} and
\ref{remark-equalities-nonconstant-rank}).
\end{example}






